% Supplement: Analysis_II_Script_v1.pdf, pp. 11--13
\section{The completion of a metric space}
\label{sec:completion_of_a_metric_space}

% This section has no counterpart in Corsin's notes. It is transcribed directly from the official
% script (S9.1.3, printed pp. 11--13), rather than enriching an existing tutor passage the way the
% rest of this document's supplements do. Included because script Exercises 9.33--9.35 walking
% through the construction are short, self-contained, and exactly the kind of material this
% chapter collects. Not included: the script's own use of the construction to build R from Q
% (S9.1.4), which it marks "extra; cf. Grundstrukturen" and which belongs to the sibling project.
% Transition: Claude Opus 5 (High)
\Cref{rem:complete_vs_closed} left one question open. We know $\mathbb{Q}$ is not complete and
$\mathbb{R}$ is, and that $\mathbb{R}$ is built to repair exactly that defect. Is the repair a
one-off trick special to the reals, or can any incomplete space be fixed the same way?

It can, and the construction below is the general answer. Every metric space embeds isometrically,
and densely, into a complete one, obtained by taking the space of its own Cauchy sequences and
declaring two of them equal when they ought to have the same limit. The reals from the rationals
is then just the first instance.

\begin{definition}[Completion of a metric space]
\label{def:completion_of_metric_space}
Let $(X,d)$ be a metric space. Write $\mathcal{C}_X$ for the set of all Cauchy sequences in $X$,
and define an equivalence relation on $\mathcal{C}_X$ by
\[(x_n)_{n=0}^\infty \sim (y_n)_{n=0}^\infty \quad :\Longleftrightarrow\quad \lim_{n\to\infty}d(x_n,y_n) = 0.\]
The quotient set $\overline X := \mathcal{C}_X/{\sim}$, equipped with the map
\[\overline d\big([(x_n)_{n=0}^\infty],\,[(y_n)_{n=0}^\infty]\big) := \lim_{n\to\infty}d(x_n,y_n),\]
is called the \newterm{completion} (\germanterm{Vervollst\"andigung}) of $(X,d)$. The injection
$\iota : X \to \overline X$ sending $x$ to the class of the constant sequence $(x,x,x,\dots)$ is
called the \newterm{canonical embedding}. For all $x,y \in X$,
\[d(x,y) = \overline d\big(\iota(x),\iota(y)\big),\]
which in particular shows $\iota$ is injective.
\end{definition}

% Supplement: Analysis_II_Script_v1.pdf, p. 11
\begin{exercise}[Well-definedness of the completion]
\label{ex:completion_well_defined}
Show that the objects introduced in \cref{def:completion_of_metric_space} are well defined. In
particular, verify that $\overline d$ is indeed a metric on $\overline X$.
\end{exercise}

% Supplement: Analysis_II_Script_v1.pdf, p. 11
\begin{exercise}[The completion is complete]
\label{ex:completion_is_complete}
As the name suggests, the completion $(\overline X,\overline d)$ of a metric space is complete,
meaning every Cauchy sequence in $\overline X$ converges. A sequence in $\overline X$ is
essentially a sequence of sequences, i.e.\
$\big[(x_{m,n})_{n=0}^\infty\big]_{m=0}^\infty$. Show that $\big[(x_{n,n})_{n=0}^\infty\big]$ is a
limit of this sequence.
\end{exercise}

% Supplement: Analysis_II_Script_v1.pdf, p. 11
\begin{exercise}[Universal property of the completion]
\label{ex:completion_universal_property}
Let $(X,d)$ be a metric space with completion $(\overline X,\overline d)$. Let $(Y,d_Y)$ be a
complete metric space, and let $f : X \to Y$ be a function such that
\[d(x,y) = d_Y\big(f(x),f(y)\big) \qquad \text{for all } x,y \in X.\]
Show that there exists a unique function $\overline f : \overline X \to Y$ such that
$f = \overline f \circ \iota_X$ and
\[\overline d(\overline x,\overline y) = d_Y\big(\overline f(\overline x),\overline f(\overline y)\big)
  \qquad \text{for all } \overline x,\overline y \in \overline X.\]
This can be interpreted as: \qt{$\overline X$ is the \emph{smallest} complete metric space
containing $X$.}
\end{exercise}

\begin{remark}
\Cref{ex:completion_universal_property} is what makes \qt{the} completion well-posed rather than
just \qt{a} completion: any other complete space that $X$ embeds into isometrically receives a
unique isometric copy of $\overline X$ as well, via $\overline f$. This is the same
uniqueness-up-to-isomorphism pattern behind \qt{the} real numbers as \qt{the} complete ordered
field, and it is exactly how the script uses this construction in $\S9.1.4$, taking
$X = \mathbb{Q}$.
\end{remark}
